\clearpage\phantomsection\label{sujet:4-1}\addcontentsline{toc}{chapter}{Test 3 · Practice 1}

\section*{Test 3 · Practice 1}

{\sffamily\bfseries\color{teal}PRACTICE PAPER · TRAINING\par}

\textbf{Duration: 2 h · Marked out of 20 · Four independent exercises.}\\Calculator allowed. Every answer must be justified; incomplete working will be taken into account. Give exact values; round probabilities to $10^{-6}$ and other approximations to $10^{-4}$, unless stated otherwise.

\smallskip\textbf{Name:}\enspace\dotfill\quad\textbf{Class:}\enspace\dotfill\quad\textbf{Date:}\enspace\dotfill

\subsection*{Exercise 1 · Limits and asymptotes \hfill 5 marks}Consider the function $f$ defined by $f(x)=\frac{5x+4}{x-1}$. Let $\mathcal C_f$ denote its graph in a coordinate plane.

\begin{enumerate}

\item Determine the domain of $f$. Find two real numbers $A$ and $B$ such that, for all $x\ne1$, $f(x)=A+\frac B{x-1}$. \hfill\textit{(1 pt)}

\item Determine the limits of $f$ at $-\infty$ and at $+\infty$. \hfill\textit{(1 pt)}

\item Determine the limits of $f$ at $1^-$ and at $1^+$. \hfill\textit{(1 pt)}

\item Give the equations of the asymptotes of $\mathcal C_f$. State the position of the curve relative to its horizontal asymptote. \hfill\textit{(1 pt)}

\item Independently of the previous questions, determine $\displaystyle\lim_{x\to+\infty}x^{5}e^{-1x}$ with justification. \hfill\textit{(1 pt)}

\end{enumerate}

\subsection*{Exercise 2 · Variations, convexity and equation \hfill 5 marks}Consider the function $f$ defined on $\R$ par $f(x)=x^3+4x-5$. Let $\mathcal C_f$ denote its graph.

\begin{enumerate}

\item Calculate $f'(x)$, then study the direction of variation of $f$ on $\R$. \hfill\textit{(1 pt)}

\item Determine the limits of $f$ at $-\infty$ and at $+\infty$. \hfill\textit{(0.5 pt)}

\item Prove that the equation $f(x)=0$ has a unique real solution $\alpha$. Give its exact value if it is an integer; otherwise, bound it between two consecutive integers. \hfill\textit{(1.5 pt)}

\item Study the convexity of $f$, then determine the coordinates of the point of inflection of $\mathcal C_f$. \hfill\textit{(1 pt)}

\item Determine an equation of the tangent to $\mathcal C_f$ at the point of abscissa $0$, then study the position of the curve relative to this tangent. \hfill\textit{(1 pt)}

\end{enumerate}

\clearpage\noindent{\sffamily\bfseries\color{navy}Test 3 · Practice 1 · continued}\par

\subsection*{Exercise 3 · An exponential curve \hfill 5 marks}Consider the function $f$ defined on $[0;+\infty[$ par $f(x)=(x+3)e^{-x}$. Let $\mathcal C_f$ denote its graph.

\begin{enumerate}

\item Calculate $f(0)$ and determine the limit of $f$ at $+\infty$. \hfill\textit{(0.5 pt)}

\item Calculate $f'(x)$, then draw the variation table of $f$. \hfill\textit{(1 pt)}

\item Calculate $f''(x)$, then study the convexity of $f$. \hfill\textit{(1 pt)}

\item Prove that the equation $f(x)=\frac{3}{2}$ has a unique solution $\alpha$ in $[0;+\infty[$. \hfill\textit{(1.5 pt)}

\item Determine an equation of the tangent to $\mathcal C_f$ at the point of abscissa $0$. Using convexity, give an inequality between $f(x)$ and the affine expression of this tangent, for all $x\ge0$. \hfill\textit{(1 pt)}

\end{enumerate}

\subsection*{Exercise 4 · Forming a group \hfill 5 marks}A group consists of $6$ teachers and $10$ students, all distinct. A team of four people is chosen; the order of selection does not matter.

\begin{enumerate}

\item How many teams of four people can be formed? \hfill\textit{(1 pt)}

\item How many of these teams contain exactly two teachers? \hfill\textit{(1 pt)}

\item How many of these teams contain at least one teacher and at least one student? \hfill\textit{(1 pt)}

\item A team containing exactly two teachers is chosen, then among its members a teacher spokesperson and a student reporter are designated. Count the choices of the team together with these two assignments. \hfill\textit{(1 pt)}

\item Explain why the number of ordered lists of four distinct people, divided by $4!$, gives the number of teams obtained in question 1. \hfill\textit{(1 pt)}

\end{enumerate}

\ifprof\clearpage

\section*{Solutions · Test 3 · Practice 1}\label{cor:4-1}

\subsection*{Exercise 1 · Limits and asymptotes}\begin{enumerate}

\item \textit{[1 pt]} The denominator vanishes only at $1$: $D_f=\R\setminus\{1\}$. Since $5x+4=5(x-1)+9$, we obtain $f(x)=5+\dfrac{9}{x-1}$.

\item \textit{[1 pt]} Both limits equal $5$.

\item \textit{[1 pt]} As $x\to 1^-$, $x-1\to0$ through negative values; since $9>0$, $\frac{9}{x-1}\to-\infty$, so $f(x)\to-\infty$. On the right, the denominator tends to $0$ through positive values: the quotient, and hence $f(x)$, tend to $+\infty$.

\item \textit{[1 pt]} $x=1$ and $y=5$; the curve is below for $x<1$ and above for $x>1$.

\item \textit{[1 pt]} The comparison of growth rates of the power and the exponential directly gives $x^5e^{-x}=\frac{x^5}{e^x}\to0$ as $x\to+\infty$.

\end{enumerate}

\subsection*{Exercise 2 · Variations, convexity and equation}\begin{enumerate}

\item \textit{[1 pt]} $f$ is differentiable on $\R$ and $f'(x)=3x^2+4$. Since $3x^2\ge0$ and $4>0$, $f'(x)>0$ for all $x$; $f$ is strictly increasing on $\R$.

\item \textit{[0.5 pt]} The dominant term is $x^3$: limits $-\infty$ and $+\infty$.

\item \textit{[1.5 pt]} $f$ is a polynomial, hence continuous on $\R$. Its strict increase and its limits $-\infty,+\infty$ allow us to apply the corollary of the intermediate value theorem (IVT): there is a unique zero $\alpha$. $f(1)=0$, so the solution is exactly $1$.

\item \textit{[1 pt]} $f''(x)=6x$: concave for $x\le0$, convex for $x\ge0$, point of inflection $(0,-5)$.

\item \textit{[1 pt]} Tangent $y=4x-5$; difference $x^3$, negative to the left and positive to the right of $0$.

\end{enumerate}

\subsection*{Exercise 3 · An exponential curve}\begin{enumerate}

\item \textit{[0.5 pt]} $f(0)=3$; limit $0$ by comparison of growth rates.

\item \textit{[1 pt]} By the product rule,
\[f'(x)=e^{-x}-(x+3)e^{-x}=(-x-2)e^{-x}.\]
The two factors are respectively strictly negative and positive on $[0,+\infty[$; hence $f'<0$.
\[\begin{array}{c|ccc}x&0&&+\infty\\\hline f'(x)&&-&\\\hline f(x)&3&\searrow&0\end{array}\]

\item \textit{[1 pt]} Differentiating the previous product,
\[f''(x)=-e^{-x}-(-x-2)e^{-x}=(x+1)e^{-x}.\]
On the domain, $x+1\ge0$ and $e^{-x}>0$: $f$ is convex. There is no change of convexity in the interior of the domain, hence no point of inflection.

\item \textit{[1.5 pt]} Continuity, strict decrease, $f(0)=3>\frac{3}{2}$ and limit $0<\frac{3}{2}$: exactly one solution.

\item \textit{[1 pt]} The tangent at $0$ has equation $y=f(0)+f'(0)x=3-2x$. A convex function lies above each of its tangents, so $(x+3)e^{-x}\ge 3-2x$ for all $x\ge0$.

\end{enumerate}

\subsection*{Exercise 4 · Forming a group}\begin{enumerate}

\item \textit{[1 pt]} Order does not distinguish teams: we must choose a subset of size $4$ from $16$ people, that is $\binom{16}4=1820$ teams.

\item \textit{[1 pt]} We choose two teachers from $6$ and two students from $10$. Each pair of choices gives a unique team: $\binom{6}2\binom{10}2=675$.

\item \textit{[1 pt]} The excluded teams are those made up only of teachers or only of students; these two categories are disjoint. The required number is $\binom{16}4-\binom{6}4-\binom{10}4=1595$.

\item \textit{[1 pt]} Each team yields $2\times2$ assignments: $2700$.

\item \textit{[1 pt]} Each team of four distinct people corresponds to exactly $4!$ ordered lists. This multiplicity is constant, so the division removes the multiple counting.

\end{enumerate}

\fi

\clearpage\phantomsection\label{sujet:4-2}\addcontentsline{toc}{chapter}{Test 3 · Practice 2}

\section*{Test 3 · Practice 2}

{\sffamily\bfseries\color{teal}PRACTICE PAPER · TRAINING\par}

\textbf{Duration: 2 h · Marked out of 20 · Four independent exercises.}\\Calculator allowed. Every answer must be justified; incomplete working will be taken into account. Give exact values; round probabilities to $10^{-6}$ and other approximations to $10^{-4}$, unless stated otherwise.

\smallskip\textbf{Name:}\enspace\dotfill\quad\textbf{Class:}\enspace\dotfill\quad\textbf{Date:}\enspace\dotfill

\subsection*{Exercise 1 · Limits and asymptotes \hfill 5 marks}Consider the function $f$ defined by $f(x)=\frac{4x+6}{x-2}$. Let $\mathcal C_f$ denote its graph in a coordinate plane.

\begin{enumerate}

\item Determine the domain of $f$. Find two real numbers $A$ and $B$ such that, for all $x\ne2$, $f(x)=A+\frac B{x-2}$. \hfill\textit{(1 pt)}

\item Determine the limits of $f$ at $-\infty$ and at $+\infty$. \hfill\textit{(1 pt)}

\item Determine the limits of $f$ at $2^-$ and at $2^+$. \hfill\textit{(1 pt)}

\item Give the equations of the asymptotes of $\mathcal C_f$. State the position of the curve relative to its horizontal asymptote. \hfill\textit{(1 pt)}

\item Independently of the previous questions, determine $\displaystyle\lim_{x\to+\infty}x^{4}e^{-2x}$ with justification. \hfill\textit{(1 pt)}

\end{enumerate}

\subsection*{Exercise 2 · Variations, convexity and equation \hfill 5 marks}Consider the function $f$ defined on $\R$ par $f(x)=x^3+3x-7$. Let $\mathcal C_f$ denote its graph.

\begin{enumerate}

\item Calculate $f'(x)$, then study the direction of variation of $f$ on $\R$. \hfill\textit{(1 pt)}

\item Determine the limits of $f$ at $-\infty$ and at $+\infty$. \hfill\textit{(0.5 pt)}

\item Prove that the equation $f(x)=0$ has a unique real solution $\alpha$. Give its exact value if it is an integer; otherwise, bound it between two consecutive integers. \hfill\textit{(1.5 pt)}

\item Study the convexity of $f$, then determine the coordinates of the point of inflection of $\mathcal C_f$. \hfill\textit{(1 pt)}

\item Determine an equation of the tangent to $\mathcal C_f$ at the point of abscissa $0$, then study the position of the curve relative to this tangent. \hfill\textit{(1 pt)}

\end{enumerate}

\clearpage\noindent{\sffamily\bfseries\color{navy}Test 3 · Practice 2 · continued}\par

\subsection*{Exercise 3 · An exponential curve \hfill 5 marks}Consider the function $f$ defined on $[0;+\infty[$ par $f(x)=(x+2)e^{-x}$. Let $\mathcal C_f$ denote its graph.

\begin{enumerate}

\item Calculate $f(0)$ and determine the limit of $f$ at $+\infty$. \hfill\textit{(0.5 pt)}

\item Calculate $f'(x)$, then draw the variation table of $f$. \hfill\textit{(1 pt)}

\item Calculate $f''(x)$, then study the convexity of $f$. \hfill\textit{(1 pt)}

\item Prove that the equation $f(x)=1$ has a unique solution $\alpha$ in $[0;+\infty[$. \hfill\textit{(1.5 pt)}

\item Determine an equation of the tangent to $\mathcal C_f$ at the point of abscissa $0$. Using convexity, give an inequality between $f(x)$ and the affine expression of this tangent, for all $x\ge0$. \hfill\textit{(1 pt)}

\end{enumerate}

\subsection*{Exercise 4 · Forming a group \hfill 5 marks}A group consists of $9$ teachers and $11$ students, all distinct. A team of four people is chosen; the order of selection does not matter.

\begin{enumerate}

\item How many teams of four people can be formed? \hfill\textit{(1 pt)}

\item How many of these teams contain exactly two teachers? \hfill\textit{(1 pt)}

\item How many of these teams contain at least one teacher and at least one student? \hfill\textit{(1 pt)}

\item A team containing exactly two teachers is chosen, then among its members a teacher spokesperson and a student reporter are designated. Count the choices of the team together with these two assignments. \hfill\textit{(1 pt)}

\item Explain why the number of ordered lists of four distinct people, divided by $4!$, gives the number of teams obtained in question 1. \hfill\textit{(1 pt)}

\end{enumerate}

\ifprof\clearpage

\section*{Solutions · Test 3 · Practice 2}\label{cor:4-2}

\subsection*{Exercise 1 · Limits and asymptotes}\begin{enumerate}

\item \textit{[1 pt]} The denominator vanishes only at $2$: $D_f=\R\setminus\{2\}$. Since $4x+6=4(x-2)+14$, we obtain $f(x)=4+\dfrac{14}{x-2}$.

\item \textit{[1 pt]} Both limits equal $4$.

\item \textit{[1 pt]} As $x\to 2^-$, $x-2\to0$ through negative values; since $14>0$, $\frac{14}{x-2}\to-\infty$, so $f(x)\to-\infty$. On the right, the denominator tends to $0$ through positive values: the quotient, and hence $f(x)$, tend to $+\infty$.

\item \textit{[1 pt]} $x=2$ and $y=4$; the curve is below for $x<2$ and above for $x>2$.

\item \textit{[1 pt]} Let $t=2x$. Then $t\to+\infty$ and $x^4e^{-2x}=\frac{1}{16}\frac{t^4}{e^t}$. Comparison of growth rates gives $\frac{t^4}{e^t}\to0$; the required limit is therefore $0$.

\end{enumerate}

\subsection*{Exercise 2 · Variations, convexity and equation}\begin{enumerate}

\item \textit{[1 pt]} $f$ is differentiable on $\R$ and $f'(x)=3x^2+3$. Since $3x^2\ge0$ and $3>0$, $f'(x)>0$ for all $x$; $f$ is strictly increasing on $\R$.

\item \textit{[0.5 pt]} The dominant term is $x^3$: limits $-\infty$ and $+\infty$.

\item \textit{[1.5 pt]} $f$ is a polynomial, hence continuous on $\R$. Its strict increase and its limits $-\infty,+\infty$ allow us to apply the corollary of the intermediate value theorem (IVT): there is a unique zero $\alpha$. $f(1)=-3<0$ and $f(2)=7>0$; so $1<\alpha<2$.

\item \textit{[1 pt]} $f''(x)=6x$: concave for $x\le0$, convex for $x\ge0$, point of inflection $(0,-7)$.

\item \textit{[1 pt]} Tangent $y=3x-7$; difference $x^3$, negative to the left and positive to the right of $0$.

\end{enumerate}

\subsection*{Exercise 3 · An exponential curve}\begin{enumerate}

\item \textit{[0.5 pt]} $f(0)=2$; limit $0$ by comparison of growth rates.

\item \textit{[1 pt]} By the product rule,
\[f'(x)=e^{-x}-(x+2)e^{-x}=(-x-1)e^{-x}.\]
The two factors are respectively strictly negative and positive on $[0,+\infty[$; hence $f'<0$.
\[\begin{array}{c|ccc}x&0&&+\infty\\\hline f'(x)&&-&\\\hline f(x)&2&\searrow&0\end{array}\]

\item \textit{[1 pt]} Differentiating the previous product,
\[f''(x)=-e^{-x}-(-x-1)e^{-x}=xe^{-x}.\]
On the domain, $x+0\ge0$ and $e^{-x}>0$: $f$ is convex. There is no change of convexity in the interior of the domain, hence no point of inflection.

\item \textit{[1.5 pt]} Continuity, strict decrease, $f(0)=2>1$ and limit $0<1$: exactly one solution.

\item \textit{[1 pt]} The tangent at $0$ has equation $y=f(0)+f'(0)x=2-x$. A convex function lies above each of its tangents, so $(x+2)e^{-x}\ge 2-x$ for all $x\ge0$.

\end{enumerate}

\subsection*{Exercise 4 · Forming a group}\begin{enumerate}

\item \textit{[1 pt]} Order does not distinguish teams: we must choose a subset of size $4$ from $20$ people, that is $\binom{20}4=4845$ teams.

\item \textit{[1 pt]} We choose two teachers from $9$ and two students from $11$. Each pair of choices gives a unique team: $\binom{9}2\binom{11}2=1980$.

\item \textit{[1 pt]} The excluded teams are those made up only of teachers or only of students; these two categories are disjoint. The required number is $\binom{20}4-\binom{9}4-\binom{11}4=4389$.

\item \textit{[1 pt]} Each team yields $2\times2$ assignments: $7920$.

\item \textit{[1 pt]} Each team of four distinct people corresponds to exactly $4!$ ordered lists. This multiplicity is constant, so the division removes the multiple counting.

\end{enumerate}

\fi

\ifprof

\clearpage\phantomsection\label{sujet:4-3}\addcontentsline{toc}{chapter}{Test 3 · A}

\section*{Test 3 · A}

{\sffamily\bfseries\color{amber}PAPER RESERVED FOR ASSESSMENT\par}

\textbf{Duration: 2 h · Marked out of 20 · Four independent exercises.}\\Calculator allowed. Every answer must be justified; incomplete working will be taken into account. Give exact values; round probabilities to $10^{-6}$ and other approximations to $10^{-4}$, unless stated otherwise.

\smallskip\textbf{Name:}\enspace\dotfill\quad\textbf{Class:}\enspace\dotfill\quad\textbf{Date:}\enspace\dotfill

\subsection*{Exercise 1 · Limits and asymptotes \hfill 5 marks}Consider the function $f$ defined by $f(x)=\frac{3x+8}{x-3}$. Let $\mathcal C_f$ denote its graph in a coordinate plane.

\begin{enumerate}

\item Determine the domain of $f$. Find two real numbers $A$ and $B$ such that, for all $x\ne3$, $f(x)=A+\frac B{x-3}$. \hfill\textit{(1 pt)}

\item Determine the limits of $f$ at $-\infty$ and at $+\infty$. \hfill\textit{(1 pt)}

\item Determine the limits of $f$ at $3^-$ and at $3^+$. \hfill\textit{(1 pt)}

\item Give the equations of the asymptotes of $\mathcal C_f$. State the position of the curve relative to its horizontal asymptote. \hfill\textit{(1 pt)}

\item Independently of the previous questions, determine $\displaystyle\lim_{x\to+\infty}x^{3}e^{-3x}$ with justification. \hfill\textit{(1 pt)}

\end{enumerate}

\subsection*{Exercise 2 · Variations, convexity and equation \hfill 5 marks}Consider the function $f$ defined on $\R$ par $f(x)=x^3+2x-4$. Let $\mathcal C_f$ denote its graph.

\begin{enumerate}

\item Calculate $f'(x)$, then study the direction of variation of $f$ on $\R$. \hfill\textit{(1 pt)}

\item Determine the limits of $f$ at $-\infty$ and at $+\infty$. \hfill\textit{(0.5 pt)}

\item Prove that the equation $f(x)=0$ has a unique real solution $\alpha$. Give its exact value if it is an integer; otherwise, bound it between two consecutive integers. \hfill\textit{(1.5 pt)}

\item Study the convexity of $f$, then determine the coordinates of the point of inflection of $\mathcal C_f$. \hfill\textit{(1 pt)}

\item Determine an equation of the tangent to $\mathcal C_f$ at the point of abscissa $0$, then study the position of the curve relative to this tangent. \hfill\textit{(1 pt)}

\end{enumerate}

\clearpage\noindent{\sffamily\bfseries\color{navy}Test 3 · A · continued}\par

\subsection*{Exercise 3 · An exponential curve \hfill 5 marks}Consider the function $f$ defined on $[0;+\infty[$ par $f(x)=(x+5)e^{-x}$. Let $\mathcal C_f$ denote its graph.

\begin{enumerate}

\item Calculate $f(0)$ and determine the limit of $f$ at $+\infty$. \hfill\textit{(0.5 pt)}

\item Calculate $f'(x)$, then draw the variation table of $f$. \hfill\textit{(1 pt)}

\item Calculate $f''(x)$, then study the convexity of $f$. \hfill\textit{(1 pt)}

\item Prove that the equation $f(x)=\frac{5}{2}$ has a unique solution $\alpha$ in $[0;+\infty[$. \hfill\textit{(1.5 pt)}

\item Determine an equation of the tangent to $\mathcal C_f$ at the point of abscissa $0$. Using convexity, give an inequality between $f(x)$ and the affine expression of this tangent, for all $x\ge0$. \hfill\textit{(1 pt)}

\end{enumerate}

\subsection*{Exercise 4 · Forming a group \hfill 5 marks}A group consists of $8$ teachers and $8$ students, all distinct. A team of four people is chosen; the order of selection does not matter.

\begin{enumerate}

\item How many teams of four people can be formed? \hfill\textit{(1 pt)}

\item How many of these teams contain exactly two teachers? \hfill\textit{(1 pt)}

\item How many of these teams contain at least one teacher and at least one student? \hfill\textit{(1 pt)}

\item A team containing exactly two teachers is chosen, then among its members a teacher spokesperson and a student reporter are designated. Count the choices of the team together with these two assignments. \hfill\textit{(1 pt)}

\item Explain why the number of ordered lists of four distinct people, divided by $4!$, gives the number of teams obtained in question 1. \hfill\textit{(1 pt)}

\end{enumerate}

\ifprof\clearpage

\section*{Solutions · Test 3 · A}\label{cor:4-3}

\subsection*{Exercise 1 · Limits and asymptotes}\begin{enumerate}

\item \textit{[1 pt]} The denominator vanishes only at $3$: $D_f=\R\setminus\{3\}$. Since $3x+8=3(x-3)+17$, we obtain $f(x)=3+\dfrac{17}{x-3}$.

\item \textit{[1 pt]} Both limits equal $3$.

\item \textit{[1 pt]} As $x\to 3^-$, $x-3\to0$ through negative values; since $17>0$, $\frac{17}{x-3}\to-\infty$, so $f(x)\to-\infty$. On the right, the denominator tends to $0$ through positive values: the quotient, and hence $f(x)$, tend to $+\infty$.

\item \textit{[1 pt]} $x=3$ and $y=3$; the curve is below for $x<3$ and above for $x>3$.

\item \textit{[1 pt]} Let $t=3x$. Then $t\to+\infty$ and $x^3e^{-3x}=\frac{1}{27}\frac{t^3}{e^t}$. Comparison of growth rates gives $\frac{t^3}{e^t}\to0$; the required limit is therefore $0$.

\end{enumerate}

\subsection*{Exercise 2 · Variations, convexity and equation}\begin{enumerate}

\item \textit{[1 pt]} $f$ is differentiable on $\R$ and $f'(x)=3x^2+2$. Since $3x^2\ge0$ and $2>0$, $f'(x)>0$ for all $x$; $f$ is strictly increasing on $\R$.

\item \textit{[0.5 pt]} The dominant term is $x^3$: limits $-\infty$ and $+\infty$.

\item \textit{[1.5 pt]} $f$ is a polynomial, hence continuous on $\R$. Its strict increase and its limits $-\infty,+\infty$ allow us to apply the corollary of the intermediate value theorem (IVT): there is a unique zero $\alpha$. $f(1)=-1<0$ and $f(2)=8>0$; so $1<\alpha<2$.

\item \textit{[1 pt]} $f''(x)=6x$: concave for $x\le0$, convex for $x\ge0$, point of inflection $(0,-4)$.

\item \textit{[1 pt]} Tangent $y=2x-4$; difference $x^3$, negative to the left and positive to the right of $0$.

\end{enumerate}

\subsection*{Exercise 3 · An exponential curve}\begin{enumerate}

\item \textit{[0.5 pt]} $f(0)=5$; limit $0$ by comparison of growth rates.

\item \textit{[1 pt]} By the product rule,
\[f'(x)=e^{-x}-(x+5)e^{-x}=(-x-4)e^{-x}.\]
The two factors are respectively strictly negative and positive on $[0,+\infty[$; hence $f'<0$.
\[\begin{array}{c|ccc}x&0&&+\infty\\\hline f'(x)&&-&\\\hline f(x)&5&\searrow&0\end{array}\]

\item \textit{[1 pt]} Differentiating the previous product,
\[f''(x)=-e^{-x}-(-x-4)e^{-x}=(x+3)e^{-x}.\]
On the domain, $x+3\ge0$ and $e^{-x}>0$: $f$ is convex. There is no change of convexity in the interior of the domain, hence no point of inflection.

\item \textit{[1.5 pt]} Continuity, strict decrease, $f(0)=5>\frac{5}{2}$ and limit $0<\frac{5}{2}$: exactly one solution.

\item \textit{[1 pt]} The tangent at $0$ has equation $y=f(0)+f'(0)x=5-4x$. A convex function lies above each of its tangents, so $(x+5)e^{-x}\ge 5-4x$ for all $x\ge0$.

\end{enumerate}

\subsection*{Exercise 4 · Forming a group}\begin{enumerate}

\item \textit{[1 pt]} Order does not distinguish teams: we must choose a subset of size $4$ from $16$ people, that is $\binom{16}4=1820$ teams.

\item \textit{[1 pt]} We choose two teachers from $8$ and two students from $8$. Each pair of choices gives a unique team: $\binom{8}2\binom{8}2=784$.

\item \textit{[1 pt]} The excluded teams are those made up only of teachers or only of students; these two categories are disjoint. The required number is $\binom{16}4-\binom{8}4-\binom{8}4=1680$.

\item \textit{[1 pt]} Each team yields $2\times2$ assignments: $3136$.

\item \textit{[1 pt]} Each team of four distinct people corresponds to exactly $4!$ ordered lists. This multiplicity is constant, so the division removes the multiple counting.

\end{enumerate}

\fi

\fi

\ifprof

\clearpage\phantomsection\label{sujet:4-4}\addcontentsline{toc}{chapter}{Test 3 · B}

\section*{Test 3 · B}

{\sffamily\bfseries\color{amber}PAPER RESERVED FOR ASSESSMENT\par}

\textbf{Duration: 2 h · Marked out of 20 · Four independent exercises.}\\Calculator allowed. Every answer must be justified; incomplete working will be taken into account. Give exact values; round probabilities to $10^{-6}$ and other approximations to $10^{-4}$, unless stated otherwise.

\smallskip\textbf{Name:}\enspace\dotfill\quad\textbf{Class:}\enspace\dotfill\quad\textbf{Date:}\enspace\dotfill

\subsection*{Exercise 1 · Limits and asymptotes \hfill 5 marks}Consider the function $f$ defined by $f(x)=\frac{2x+5}{x-1}$. Let $\mathcal C_f$ denote its graph in a coordinate plane.

\begin{enumerate}

\item Determine the domain of $f$. Find two real numbers $A$ and $B$ such that, for all $x\ne1$, $f(x)=A+\frac B{x-1}$. \hfill\textit{(1 pt)}

\item Determine the limits of $f$ at $-\infty$ and at $+\infty$. \hfill\textit{(1 pt)}

\item Determine the limits of $f$ at $1^-$ and at $1^+$. \hfill\textit{(1 pt)}

\item Give the equations of the asymptotes of $\mathcal C_f$. State the position of the curve relative to its horizontal asymptote. \hfill\textit{(1 pt)}

\item Independently of the previous questions, determine $\displaystyle\lim_{x\to+\infty}x^{2}e^{-1x}$ with justification. \hfill\textit{(1 pt)}

\end{enumerate}

\subsection*{Exercise 2 · Variations, convexity and equation \hfill 5 marks}Consider the function $f$ defined on $\R$ par $f(x)=x^3+5x-5$. Let $\mathcal C_f$ denote its graph.

\begin{enumerate}

\item Calculate $f'(x)$, then study the direction of variation of $f$ on $\R$. \hfill\textit{(1 pt)}

\item Determine the limits of $f$ at $-\infty$ and at $+\infty$. \hfill\textit{(0.5 pt)}

\item Prove that the equation $f(x)=0$ has a unique real solution $\alpha$. Give its exact value if it is an integer; otherwise, bound it between two consecutive integers. \hfill\textit{(1.5 pt)}

\item Study the convexity of $f$, then determine the coordinates of the point of inflection of $\mathcal C_f$. \hfill\textit{(1 pt)}

\item Determine an equation of the tangent to $\mathcal C_f$ at the point of abscissa $0$, then study the position of the curve relative to this tangent. \hfill\textit{(1 pt)}

\end{enumerate}

\clearpage\noindent{\sffamily\bfseries\color{navy}Test 3 · B · continued}\par

\subsection*{Exercise 3 · An exponential curve \hfill 5 marks}Consider the function $f$ defined on $[0;+\infty[$ par $f(x)=(x+4)e^{-x}$. Let $\mathcal C_f$ denote its graph.

\begin{enumerate}

\item Calculate $f(0)$ and determine the limit of $f$ at $+\infty$. \hfill\textit{(0.5 pt)}

\item Calculate $f'(x)$, then draw the variation table of $f$. \hfill\textit{(1 pt)}

\item Calculate $f''(x)$, then study the convexity of $f$. \hfill\textit{(1 pt)}

\item Prove that the equation $f(x)=2$ has a unique solution $\alpha$ in $[0;+\infty[$. \hfill\textit{(1.5 pt)}

\item Determine an equation of the tangent to $\mathcal C_f$ at the point of abscissa $0$. Using convexity, give an inequality between $f(x)$ and the affine expression of this tangent, for all $x\ge0$. \hfill\textit{(1 pt)}

\end{enumerate}

\subsection*{Exercise 4 · Forming a group \hfill 5 marks}A group consists of $7$ teachers and $10$ students, all distinct. A team of four people is chosen; the order of selection does not matter.

\begin{enumerate}

\item How many teams of four people can be formed? \hfill\textit{(1 pt)}

\item How many of these teams contain exactly two teachers? \hfill\textit{(1 pt)}

\item How many of these teams contain at least one teacher and at least one student? \hfill\textit{(1 pt)}

\item A team containing exactly two teachers is chosen, then among its members a teacher spokesperson and a student reporter are designated. Count the choices of the team together with these two assignments. \hfill\textit{(1 pt)}

\item Explain why the number of ordered lists of four distinct people, divided by $4!$, gives the number of teams obtained in question 1. \hfill\textit{(1 pt)}

\end{enumerate}

\ifprof\clearpage

\section*{Solutions · Test 3 · B}\label{cor:4-4}

\subsection*{Exercise 1 · Limits and asymptotes}\begin{enumerate}

\item \textit{[1 pt]} The denominator vanishes only at $1$: $D_f=\R\setminus\{1\}$. Since $2x+5=2(x-1)+7$, we obtain $f(x)=2+\dfrac{7}{x-1}$.

\item \textit{[1 pt]} Both limits equal $2$.

\item \textit{[1 pt]} As $x\to 1^-$, $x-1\to0$ through negative values; since $7>0$, $\frac{7}{x-1}\to-\infty$, so $f(x)\to-\infty$. On the right, the denominator tends to $0$ through positive values: the quotient, and hence $f(x)$, tend to $+\infty$.

\item \textit{[1 pt]} $x=1$ and $y=2$; the curve is below for $x<1$ and above for $x>1$.

\item \textit{[1 pt]} The comparison of growth rates of the power and the exponential directly gives $x^2e^{-x}=\frac{x^2}{e^x}\to0$ as $x\to+\infty$.

\end{enumerate}

\subsection*{Exercise 2 · Variations, convexity and equation}\begin{enumerate}

\item \textit{[1 pt]} $f$ is differentiable on $\R$ and $f'(x)=3x^2+5$. Since $3x^2\ge0$ and $5>0$, $f'(x)>0$ for all $x$; $f$ is strictly increasing on $\R$.

\item \textit{[0.5 pt]} The dominant term is $x^3$: limits $-\infty$ and $+\infty$.

\item \textit{[1.5 pt]} $f$ is a polynomial, hence continuous on $\R$. Its strict increase and its limits $-\infty,+\infty$ allow us to apply the corollary of the intermediate value theorem (IVT): there is a unique zero $\alpha$. $f(0)=-5<0$ and $f(1)=1>0$; so $0<\alpha<1$.

\item \textit{[1 pt]} $f''(x)=6x$: concave for $x\le0$, convex for $x\ge0$, point of inflection $(0,-5)$.

\item \textit{[1 pt]} Tangent $y=5x-5$; difference $x^3$, negative to the left and positive to the right of $0$.

\end{enumerate}

\subsection*{Exercise 3 · An exponential curve}\begin{enumerate}

\item \textit{[0.5 pt]} $f(0)=4$; limit $0$ by comparison of growth rates.

\item \textit{[1 pt]} By the product rule,
\[f'(x)=e^{-x}-(x+4)e^{-x}=(-x-3)e^{-x}.\]
The two factors are respectively strictly negative and positive on $[0,+\infty[$; hence $f'<0$.
\[\begin{array}{c|ccc}x&0&&+\infty\\\hline f'(x)&&-&\\\hline f(x)&4&\searrow&0\end{array}\]

\item \textit{[1 pt]} Differentiating the previous product,
\[f''(x)=-e^{-x}-(-x-3)e^{-x}=(x+2)e^{-x}.\]
On the domain, $x+2\ge0$ and $e^{-x}>0$: $f$ is convex. There is no change of convexity in the interior of the domain, hence no point of inflection.

\item \textit{[1.5 pt]} Continuity, strict decrease, $f(0)=4>2$ and limit $0<2$: exactly one solution.

\item \textit{[1 pt]} The tangent at $0$ has equation $y=f(0)+f'(0)x=4-3x$. A convex function lies above each of its tangents, so $(x+4)e^{-x}\ge 4-3x$ for all $x\ge0$.

\end{enumerate}

\subsection*{Exercise 4 · Forming a group}\begin{enumerate}

\item \textit{[1 pt]} Order does not distinguish teams: we must choose a subset of size $4$ from $17$ people, that is $\binom{17}4=2380$ teams.

\item \textit{[1 pt]} We choose two teachers from $7$ and two students from $10$. Each pair of choices gives a unique team: $\binom{7}2\binom{10}2=945$.

\item \textit{[1 pt]} The excluded teams are those made up only of teachers or only of students; these two categories are disjoint. The required number is $\binom{17}4-\binom{7}4-\binom{10}4=2135$.

\item \textit{[1 pt]} Each team yields $2\times2$ assignments: $3780$.

\item \textit{[1 pt]} Each team of four distinct people corresponds to exactly $4!$ ordered lists. This multiplicity is constant, so the division removes the multiple counting.

\end{enumerate}

\fi

\fi

